\documentclass[11pt]{article}
\usepackage[margin=1in]{geometry}
\usepackage[T1]{fontenc}
\usepackage{graphicx}
\usepackage{hyperref}
\usepackage{enumitem}
\setlength{\parskip}{0.5em}
\setlength{\parindent}{0pt}

\title{large language model agents and tool use}
\author{arXiv GraphRAG --- phase-6 deep research}
\date{\today}

\begin{document}
\maketitle

\begin{abstract}
(judge output unparseable)
\end{abstract}

\section{From Textual Prediction to Functional Agency: A Review of Large Language Model Agents and Tool Integration}

The rapid advancement of artificial intelligence has been defined by the transition of Large Language Models (LLMs) from mere statistical text predictors to sophisticated cognitive engines. While the \textit{Foundations of Large Language Models} established a robust baseline for linguistic understanding, these models initially operated within a closed semantic loop, constrained by the static knowledge contained within their training corpora. This limitation creates a significant gap between linguistic fluency and functional utility. To bridge this gap, the focus of research has shifted toward agency—the ability of a model to not only reason about a problem but to act upon the world. The importance of this evolution cannot be overstated, as the integration of tool use allows LLMs to overcome inherent limitations in mathematical reasoning, real-time information retrieval, and complex task execution, transforming them from passive responders into active participants in digital and physical environments.

This review traces the developmental arc of this transformation, beginning with the fundamental \textit{Foundations of Large Language Models} and the subsequent evolution of \textit{Reasoning and Planning Paradigms} that allow for structured, multi-step thought processes. We then investigate the \textit{Emergence of Tool-Augmented LLMs}, which represents the critical juncture where models began to interact with external APIs and software. The narrative progresses into the orchestration of \textit{Autonomous Agents and Multi-Agent Systems}, exploring how individual agentic capabilities scale into collaborative, self-correcting ecosystems. After reviewing the current landscape of \textit{Specialized Agentic Applications and Frameworks}, the review concludes at \textit{The Frontier: Advanced Agentic Reasoning and Tool Integration}, where current research efforts aim to synthesize high-order cognitive reasoning with seamless, intuitive tool manipulation to achieve true functional autonomy.

\subsection{Foundations of Large Language Models}

The architectural foundation of modern large language models (LLMs) is the Transformer, which utilizes attention mechanisms to dispense with the need for recurrence and convolutions [\texttt{arXiv:1706.03762}]. Early applications of this architecture focused on machine translation, achieving 28.4 BLEU on WMT 2014 English-to-German and 41.8 BLEU on English-to-French [\texttt{arXiv:1706.03762}]. Subsequent research diverged in how these models process context: while the Transformer provides a versatile base, BERT introduced deep bidirectional Transformers that jointly condition on both left and right context [\texttt{arXiv:1810.04805}]. This bidirectional approach yielded high performance on benchmarks such as GLUE (80.5\%) and SQuAD v1.1 (93.2\% Test F1) [\texttt{arXiv:1810.04805}]. In contrast, GPT-3 shifted the focus toward massive-scale autoregressive modeling, utilizing 175 billion parameters to enable few-shot learning through text interaction without requiring gradient updates [\texttt{arXiv:2005.14165}].

As model scale increased, the nature of model capabilities became a central point of investigation. The MMLU benchmark highlighted a significant gap between scale and reliability, noting that models often exhibit lopsided performance and a frequent inability to recognize when they are wrong, sometimes reaching only near-random accuracy on socially important subjects [\texttt{arXiv:2009.03300}]. This unpredictability is characterized as the "emergent abilities" phenomenon, where specific capabilities appear only at certain scales and cannot be predicted by extrapolating from smaller models [\texttt{arXiv:2206.07682}]. For example, the 540-billion parameter PaLM model [\texttt{arXiv:2204.02311}] demonstrates advanced multi-step reasoning and few-shot learning [\texttt{arXiv:2204.02311}].

A critical breakthrough in unlocking these capabilities is Chain-of-Thought (CoT) prompting, which elicits intermediate reasoning steps via few-shot exemplars [\texttt{arXiv:2201.11903}]. While CoT is a powerful mechanism that allows models to tackle complex arithmetic and symbolic reasoning—with PaLM 540B achieving state-of-the-art results on the GSM8K benchmark—it possesses notable weaknesses [\texttt{arXiv:2201.11903}]. Specifically, CoT relies on the manual composition of exemplars and lacks a guarantee of logical correctness, as models may generate reasoning paths that are "fluent but illogical" [\texttt{arXiv:2201.11903}].

The evolution of the field is currently moving from static language modeling toward agentic task execution. While GPT-4 demonstrates "sparks" of artificial general intelligence across domains like coding and medicine [\texttt{arXiv:2303.12712}], researchers suggest a need for paradigms beyond simple next-word prediction [\texttt{arXiv:2303.12712}]. This is addressed by emerging LLM agents that combine CoT with strategies such as Recursively Criticizes and Improves (RCI) to automate computer tasks [\texttt{arXiv:2303.17491}]. Unlike earlier approaches that required impractical amounts of expert demonstrations or task-specific reward functions, these agentic frameworks leverage the model's internal reasoning to solve complex, natural language commands [\texttt{arXiv:2303.17491}].

\subsection{Reasoning and Planning Paradigms}

The evolution of reasoning paradigms in large language models (LLMs) has transitioned from simple sequential generation to complex, multi-path, and interactive frameworks. Early approaches focused on improving the reliability of single-path reasoning by leveraging multiple candidates. For instance, \texttt{arXiv:2110.14168} addressed the failure of large transformers in multi-step mathematical reasoning by training verifiers to rank and select correct completions from multiple candidate solutions in grade school math problems. This concept of utilizing multiple paths was further refined by \texttt{arXiv:2203.11171}, which introduced a self-consistency decoding strategy. Rather than relying on a single verifier, this method samples diverse reasoning paths and selects the most consistent answer via marginalization, yielding significant improvements across benchmarks such as GSM8K (+17.9\%) and SVAMP (+11.0\%).

While these methods improve accuracy through consensus, they remain largely internal to the model's generative process. To address the limitations of standalone Chain-of-Thought (CoT) reasoning—specifically hallucination and error propagation—\texttt{arXiv:2210.03629} proposed the ReAct framework. ReAct shifts the paradigm from pure reasoning to a "reason-to-act" and "act-to-reason" synergy, interleaving verbal reasoning traces with task-specific actions using external tools like the Wikipedia API. This closed-loop architecture provides grounding in external environments, outperforming imitation methods by 34\% in ALFWorld. However, the framework's reliance on external tool quality is a notable weakness; for example, 23\% of errors in HotpotQA stemmed from non-informative search results, and the model remains susceptible to repetitive reasoning loops.

More sophisticated structural paradigms have emerged to handle complex decision-making. \texttt{arXiv:2305.10601} introduced the Tree of Thoughts (ToT) framework, which generalizes CoT by enabling exploration over "thoughts" as intermediate units. Unlike the left-to-right, token-level decisions that lack strategic lookahead, ToT allows for deliberate decision-making, including self-evaluation, looking ahead, and backtracking. This is evidenced by its performance in the Game of 24, where ToT achieved a 74\% success rate compared to only 4\% for GPT-4 with CoT.

Recent research has further specialized these architectures for agentic task execution. \texttt{arXiv:2308.03427} introduced the TPTU framework, which bifurcates agent architecture into one-step and sequential agents to handle complex real-world task planning and tool usage. Similarly, \texttt{arXiv:2310.03710} focuses on an autonomous agent designed to instruct the reasoning process of LLMs to enhance zero-shot capabilities across 29 datasets. Collectively, these developments represent a move from static, linear reasoning toward dynamic, structured, and environmentally-grounded agentic planning.

\subsection{Emergence of Tool-Augmented LLMs}

The evolution of Large Language Models (LLMs) has transitioned from purely parametric generation toward architectures that integrate external tools to overcome inherent limitations in knowledge access and logical reasoning. Early efforts focused on augmenting models with non-parametric memory to address the inability of LLMs to precisely manipulate or update world knowledge [\texttt{arXiv:2005.11401}]. Specifically, Retrieval-Augmented Generation (RAG) introduced a differentiable access mechanism to a dense vector index of Wikipedia, allowing models to condition generation on retrieved passages [\texttt{arXiv:2005.11401}]. This paradigm shifted the focus from internalizing all information to retrieving relevant external context, a concept further expanded by WebGPT, which utilized a text-based browser to navigate the web for long-form question-answering [\texttt{arXiv:2112.09332}]. While WebGPT improved performance through imitation learning and human preference modeling—outperforming human demonstrators in certain benchmarks [\texttt{arXiv:2112.09332}]—it represents a specific application of web-navigation rather than a generalized tool-use framework.

As research progressed, the focus shifted from simple information retrieval to complex task execution and symbolic reasoning. For instance, to mitigate the tendency of LLMs to make arithmetic or logical errors during step-by-step decomposition [\texttt{arXiv:2211.10435}], Program-aided Language Models (PAL) offloaded reasoning to a symbolic Python interpreter [\texttt{arXiv:2211.10435}]. This contrasts with the approach of embodied agents, such as those using zero-shot planners to decompose high-level tasks into mid-level plans [\texttt{arXiv:2201.07207}]. However, a significant tension exists in embodied settings: naive LLM plans often fail to map to admissible actions, creating a critical trade-off between the correctness of a plan and its actual executability [\texttt{arXiv:2201.07207}].

A major breakthrough occurred with the move toward self-supervised tool learning. Toolformer demonstrated that models could teach themselves to decide when and how to call APIs, integrating results directly into token prediction [\texttt{arXiv:2302.04761}]. This represents a more autonomous form of interaction compared to the manual program generation seen in PAL [\texttt{arXiv:2211.10435}]. Building on these foundations, recent frameworks have sought to scale these capabilities to massive, real-world environments. While earlier works lacked systematic evaluation of tool-use [\texttt{arXiv:2304.08354}], ToolLLM addressed this by introducing the ToolBench dataset and a neural API retriever capable of navigating over 16,000 RESTful APIs [\texttt{arXiv:2307.16789}]. ToolLLM further refined the reasoning process by employing a depth-first search-based decision tree to evaluate multiple reasoning traces [\texttt{arXiv:2307.16789}]. This progression highlights a clear trajectory: moving from static retrieval [\texttt{arXiv:2005.11401}] to specialized symbolic execution [\texttt{arXiv:2211.10435}], and finally to generalized, autonomous agents capable of mastering vast, real-world API ecosystems [\texttt{arXiv:2307.16789}].

\subsection{Autonomous Agents and Multi-Agent Systems}

The evolution of large language model (LLM) agents has transitioned from single-agent architectures to complex multi-agent systems (MAS) capable of collaborative reasoning and specialized task execution. Early frameworks like CAMEL introduced a role-playing communicative agent framework using inception prompting to guide agents toward task completion [\texttt{arXiv:2303.17760}]. However, this approach faced challenges regarding its reliance on human input and the difficulty of achieving true autonomous cooperation [\texttt{arXiv:2303.17760}]. In contrast, the Generative Agents architecture moves toward more sophisticated autonomy by extending LLMs with a three-part reasoning strategy—observation, planning, and reflection—allowing agents to store, synthesize, and dynamically retrieve memories in immersive environments [\texttt{arXiv:2304.03442}].

A significant portion of recent research focuses on optimizing MAS for specialized domains, particularly software engineering. ChatDev utilizes a chat-powered framework where specialized agents operate in a chat chain to handle design, coding, and testing [\texttt{arXiv:2307.07924}]. While effective, this model can suffer from technical inconsistencies across the waterfall development phases [\texttt{arXiv:2307.07924}]. To mitigate such issues, MetaGPT proposes an assembly line paradigm that assigns roles based on Standardized Operating Procedures (SOPs) [\texttt{arXiv:2308.00352}]. By encoding SOPs into prompt sequences, MetaGPT allows for the verification of intermediate results, addressing the logic inconsistencies and cascading hallucinations that often plague naively chained LLMs [\texttt{arXiv:2308.00352}].

Beyond structured workflows, research has explored using agent interaction to improve cognitive accuracy. The "society of minds" approach in multiagent debate involves multiple LLM instances proposing and debating responses over several rounds to improve mathematical and factual reasoning [\texttt{arXiv:2305.14325}]. This contrasts with more flexible, programmable frameworks like AutoGen, which enables customizable, conversable agents that can integrate LLMs, human inputs, and external tools [\texttt{arXiv:2308.08155}]. AutoGen's versatility allows it to span diverse domains from operations research to entertainment [\texttt{arXiv:2308.08155}]. Similarly, modularized frameworks like Agents provide fine-grained symbolic control and tool usage to solve tasks via natural language interfaces [\texttt{arXiv:2309.07870}].

Despite these advancements, the field faces systemic hurdles. Survey literature identifies a lack of a general, powerful model capable of adapting to diverse scenarios as a primary bottleneck [\texttt{arXiv:2309.07864}]. Current MAS research continues to grapple with the complexities of planning, decision-making, and the dynamic nature of agent communication [\texttt{arXiv:2402.01680}]. While frameworks have moved from simple role-playing [\texttt{arXiv:2303.17760}] to sophisticated memory-driven simulations [\texttt{arXiv:2304.03442}] and SOP-driven engineering [\texttt{arXiv:2308.00352}], achieving seamless, hallucination-free autonomy remains a central challenge.

\subsection{Specialized Agentic Applications and Frameworks}

Recent research has shifted from general-purpose language modeling toward specialized agentic frameworks designed to master complex tool interactions and embodied exploration. A primary focus in this domain is the mastery of external APIs. Gorilla addresses the tendency of LLMs to hallucinate incorrect API usage or struggle with frequently updated documentation by combining a fine-tuned LLaMA-based model with a document retriever to generate accurate calls for HuggingFace, TorchHub, and TensorHub APIs (\texttt{arXiv:2305.15334}). Expanding upon this necessity for scale, ToolLLM introduces a framework to manage over 16,000 real-world RESTful APIs across 49 categories (\texttt{arXiv:2307.16789}). While Gorilla relies on a document retriever, ToolLLM utilizes a neural API retriever and a depth-first search-based decision tree to expand the reasoning search space and evaluate multiple solution paths (\texttt{arXiv:2307.16789}). This comparison highlights a progression from targeted API generation to large-scale, multi-step reasoning across vast API ecosystems.

Beyond static API calls, agentic research has explored embodied, open-ended learning in dynamic environments. Voyager represents a departure from API-centric models by focusing on lifelong learning within Minecraft (\texttt{arXiv:2305.16291}). Unlike the retrieval-heavy architectures of Gorilla or ToolLLM, Voyager employs a three-component system: an automatic curriculum, an iterative prompting mechanism that incorporates environmental feedback, and an "ever-growing skill library" of executable code. This library allows the agent to store and retrieve complex behaviors via blackbox queries to GPT-4, resulting in significant performance gains, including 15.3x faster unlocking of tech tree milestones compared to prior state-of-the-art methods (\texttt{arXiv:2305.16291}).

The evaluation of these specialized capabilities remains a critical challenge. AgentBench provides a comprehensive benchmark consisting of eight distinct environments to assess reasoning and decision-making (\texttt{arXiv:2308.03688}). It identifies a significant performance disparity between commercial LLMs and open-source competitors (≤70B), noting deficiencies in long-term reasoning and instruction following (\texttt{arXiv:2308.03688}). This need for rigorous assessment is echoed in broader surveys, which categorize agent architectures into planning, memory, and tool-use components (\texttt{arXiv:2309.14365}) and identify persistent hurdles such as multimodality, human value alignment, and hallucinations (\texttt{arXiv:2404.04442}). Furthermore, as agents move into specialized sectors like software engineering, the ability to perceive and utilize external resources becomes the central metric for success (\texttt{arXiv:2409.02977}). Collectively, these works demonstrate a transition from models that merely "know" information to agents that can "act" through structured tool manipulation and iterative environmental interaction.

\subsection{The Frontier: Advanced Agentic Reasoning and Tool Integration}

Recent advancements in large language model (LLM) agents signal a shift from simple tool-calling to sophisticated agentic reasoning frameworks. While early architectures focused on modular components like perception, memory, and execution systems [\texttt{arXiv:2510.09244}], contemporary research emphasizes the cognitive processes driving these interactions. A significant conceptual advancement is the three-layered framework proposed by \texttt{arXiv:2601.12538}, which distinguishes between foundational single-agent reasoning, self-evolving reasoning via feedback, and collective multi-agent coordination. This work provides a rigorous mathematical formalization of the "think-act" structure using a POMDP framework, factorizing policy into internal reasoning (\$\textbackslash{}pi\_\{reason\}\$) and external execution (\$\textbackslash{}pi\_\{exec\}\$). However, the distinction between "in-context" and "post-training" methods in this taxonomy is occasionally blurred, as "Self-evolving" layers overlap with "Parametric Adaptation" mechanisms [\texttt{arXiv:2601.12538}].

The integration of tools has evolved from basic API prompting to complex, autonomous decision-making. \texttt{arXiv:2604.00835} identifies three distinct paradigms for tool use: plug-and-play prompting, supervised tool learning, and reward-driven policy learning. This evolution is exemplified by the ARTIST framework [\texttt{arXiv:2505.01441}], which utilizes outcome-based reinforcement learning (RL) to allow agents to autonomously decide when and how to invoke tools within multi-turn reasoning chains, specifically targeting mathematical reasoning and multi-turn function calling. Similarly, \texttt{arXiv:2502.04644} enhances reasoning coherence in long chains by integrating a "Mind-Map" agent to maintain a structured knowledge graph of logical relationships alongside web-search and code-execution tools.

As agents move toward higher autonomy, the focus is shifting from linear reasoning to native inference-time reasoning and open-standard tool integration, such as the Model Context Protocol (MCP) [\texttt{arXiv:2601.12560}]. This transition is reflected in the diverse application domains identified across the literature, ranging from scientific discovery and healthcare to software engineering and finance [\texttt{arXiv:2504.19678}, \texttt{arXiv:2508.17692}, \texttt{arXiv:2601.12560}].

Despite these strides, significant challenges persist. The landscape remains fragmented [\texttt{arXiv:2504.19678}], with a lack of unified taxonomies and concerns regarding security vulnerabilities in agent protocols [\texttt{arXiv:2504.19678}]. Technical hurdles include maintaining coherence in long-horizon interactions [\texttt{arXiv:2502.04644}], addressing hallucinations in action, and preventing infinite loops [\texttt{arXiv:2601.12560}]. Furthermore, the transition to real-world deployment necessitates solving complex issues of safety, alignment, verifiable planning, and scalable multi-agent coordination [\texttt{arXiv:2601.02749}].

\subsection{Tensions and Open Problems}

A primary tension in the development of autonomous agents concerns the mechanism of reasoning: whether intelligence is best elicited through prompt-based heuristics or through structured, iterative feedback loops. On one hand, early approaches focused on augmenting the model's internal logic through prompting techniques, such as the chain-of-thought reasoning demonstrated in \texttt{arXiv:2201.11903}, or the self-consistency sampling methods proposed in \texttt{arXiv:2203.11171}. However, recent critiques suggest these methods lack the robustness required for complex tasks. For instance, \texttt{arXiv:2205.11916} argues that zero-shot reasoning can bypass some of these prompting requirements, yet more recent work in \texttt{arXiv:2302.01560} suggests that simple prompting is insufficient for interactive planning, arguing instead for more sophisticated "describe, explain, plan, and select" frameworks. This debate is further complicated by the emergence of verbal reinforcement learning, where \texttt{arXiv:2303.11366} proposes using linguistic feedback to refine agent behavior, a method that has been critiqued by \texttt{arXiv:2302.01560} for failing to provide the necessary grounding for complex, multi-step execution.

Another significant disagreement involves the integration of external tools and the degree of autonomy required for effective agency. Some frameworks emphasize the importance of teaching models to call APIs through self-supervised learning, as seen in the tool-use capabilities of \texttt{arXiv:2302.04761}. This stands in contrast to earlier, more human-in-the-loop or browser-centric approaches like \texttt{arXiv:2112.09332}, which rely heavily on human feedback for web-based tasks. While \texttt{arXiv:2309.07870} argues that the self-teaching approach of \texttt{arXiv:2302.04761} is a step toward true autonomy, it also notes that such methods can struggle with the scale and complexity of real-world environments. This tension is echoed in the debate over tool mastery; while \texttt{arXiv:2210.03629} introduced the ReAct paradigm to synergize reasoning and acting, more recent large-scale tool integration efforts, such as those in \texttt{arXiv:2307.16789}, suggest that the fundamental reasoning chains used in earlier models are often inadequate for mastering the vast, diverse toolsets required in modern ecosystems.

Finally, there is a growing debate regarding the generalizability and emergent properties of these agents. While some researchers argue that scaling parameters leads to the emergence of complex new abilities, as discussed in \texttt{arXiv:2206.07682} and \texttt{arXiv:2204.02311}, others question whether these "emergent" capabilities are truly robust or merely artifacts of specific evaluation metrics. This skepticism is central to the critique in \texttt{arXiv:2509.16330}, which argues that many current agentic frameworks—including the highly-cited survey-based methodologies in \texttt{arXiv:2309.07864} and \texttt{arXiv:2503.21460}—lack the necessary generalizability to perform reliably across diverse, out-of-distribution environments. This raises the fundamental question of whether we are building truly autonomous agents or merely sophisticated pattern matchers that excel only within the narrow bounds of their training and prompting distributions.

\subsection{Outlook}

The field is currently transitioning from "prompt-centric" reasoning toward "agentic reasoning," where the focus shifts from how a model thinks to how a system acts. As highlighted by the most recent frontiers in \texttt{arXiv:2601.12538} and \texttt{arXiv:2601.12560}, the next generation of AI will likely be defined by integrated architectures that treat tool use, environmental perception, and self-correction as unified cognitive processes. The path ahead, as outlined in \texttt{arXiv:2601.02749}, lies in moving beyond static task completion toward truly autonomous, lifelong learning agents capable of navigating open-ended, real-world complexities with minimal human intervention.

\section*{References}
\begin{enumerate}\small
  \item Attention Is All You Need (2017). \texttt{arXiv:1706.03762}
  \item Proximal Policy Optimization Algorithms (2017). \texttt{arXiv:1707.06347}
  \item BERT: Pre-training of Deep Bidirectional Transformers for Language
  Understanding (2018). \texttt{arXiv:1810.04805}
  \item HotpotQA: A Dataset for Diverse, Explainable Multi-hop Question
  Answering (2018). \texttt{arXiv:1809.09600}
  \item Language Models are Few-Shot Learners (2020). \texttt{arXiv:2005.14165}
  \item Retrieval-Augmented Generation for Knowledge-Intensive NLP Tasks (2020). \texttt{arXiv:2005.11401}
  \item Measuring Massive Multitask Language Understanding (2020). \texttt{arXiv:2009.03300}
  \item ALFWorld: Aligning Text and Embodied Environments for Interactive
  Learning (2020). \texttt{arXiv:2010.03768}
  \item Evaluating Large Language Models Trained on Code (2021). \texttt{arXiv:2107.03374}
  \item Training Verifiers to Solve Math Word Problems (2021). \texttt{arXiv:2110.14168}
  \item Finetuned Language Models Are Zero-Shot Learners (2021). \texttt{arXiv:2109.01652}
  \item WebGPT: Browser-assisted question-answering with human feedback (2021). \texttt{arXiv:2112.09332}
  \item Measuring Mathematical Problem Solving With the MATH Dataset (2021). \texttt{arXiv:2103.03874}
  \item Chain-of-Thought Prompting Elicits Reasoning in Large Language Models (2022). \texttt{arXiv:2201.11903}
  \item Training language models to follow instructions with human feedback (2022). \texttt{arXiv:2203.02155}
  \item PaLM: Scaling Language Modeling with Pathways (2022). \texttt{arXiv:2204.02311}
  \item ReAct: Synergizing Reasoning and Acting in Language Models (2022). \texttt{arXiv:2210.03629}
  \item Do As I Can, Not As I Say: Grounding Language in Robotic Affordances (2022). \texttt{arXiv:2204.01691}
  \item Self-Consistency Improves Chain of Thought Reasoning in Language Models (2022). \texttt{arXiv:2203.11171}
  \item Large Language Models are Zero-Shot Reasoners (2022). \texttt{arXiv:2205.11916}
  \item Language Models as Zero-Shot Planners: Extracting Actionable Knowledge
  for Embodied Agents (2022). \texttt{arXiv:2201.07207}
  \item Inner Monologue: Embodied Reasoning through Planning with Language
  Models (2022). \texttt{arXiv:2207.05608}
  \item Emergent Abilities of Large Language Models (2022). \texttt{arXiv:2206.07682}
  \item PAL: Program-aided Language Models (2022). \texttt{arXiv:2211.10435}
  \item WebShop: Towards Scalable Real-World Web Interaction with Grounded
  Language Agents (2022). \texttt{arXiv:2207.01206}
  \item Self-Instruct: Aligning Language Models with Self-Generated Instructions (2022). \texttt{arXiv:2212.10560}
  \item An In-depth Survey of Large Language Model-based Artificial Intelligence
  Agents (2023). \texttt{arXiv:2309.14365}
  \item A Survey on Large Language Model based Autonomous Agents (2023). \texttt{arXiv:2308.11432}
  \item Large Language Models as Tool Makers (2023). \texttt{arXiv:2305.17126}
  \item TPTU: Large Language Model-based AI Agents for Task Planning and Tool Usage (2023). \texttt{arXiv:2308.03427}
  \item Agents: An Open-source Framework for Autonomous Language Agents (2023). \texttt{arXiv:2309.07870}
  \item ART: Automatic multi-step reasoning and tool-use for large language
  models (2023). \texttt{arXiv:2303.09014}
  \item ModelScope-Agent: Building Your Customizable Agent System with
  Open-source Large Language Models (2023). \texttt{arXiv:2309.00986}
  \item Agent Instructs Large Language Models to be General Zero-Shot Reasoners (2023). \texttt{arXiv:2310.03710}
  \item The Rise and Potential of Large Language Model Based Agents: A Survey (2023). \texttt{arXiv:2309.07864}
  \item Toolformer: Language Models Can Teach Themselves to Use Tools (2023). \texttt{arXiv:2302.04761}
  \item GPT-4 Technical Report (2023). \texttt{arXiv:2303.08774}
  \item LLaMA: Open and Efficient Foundation Language Models (2023). \texttt{arXiv:2302.13971}
  \item Generative Agents: Interactive Simulacra of Human Behavior (2023). \texttt{arXiv:2304.03442}
  \item Reflexion: Language Agents with Verbal Reinforcement Learning (2023). \texttt{arXiv:2303.11366}
  \item Tree of Thoughts: Deliberate Problem Solving with Large Language Models (2023). \texttt{arXiv:2305.10601}
  \item Sparks of Artificial General Intelligence: Early experiments with GPT-4 (2023). \texttt{arXiv:2303.12712}
  \item HuggingGPT: Solving AI Tasks with ChatGPT and its Friends in Hugging
  Face (2023). \texttt{arXiv:2303.17580}
  \item ToolLLM: Facilitating Large Language Models to Master 16000+ Real-world
  APIs (2023). \texttt{arXiv:2307.16789}
  \item Self-Refine: Iterative Refinement with Self-Feedback (2023). \texttt{arXiv:2303.17651}
  \item CAMEL: Communicative Agents for "Mind" Exploration of Large Language
  Model Society (2023). \texttt{arXiv:2303.17760}
  \item Gorilla: Large Language Model Connected with Massive APIs (2023). \texttt{arXiv:2305.15334}
  \item MetaGPT: Meta Programming for A Multi-Agent Collaborative Framework (2023). \texttt{arXiv:2308.00352}
  \item Voyager: An Open-Ended Embodied Agent with Large Language Models (2023). \texttt{arXiv:2305.16291}
  \item Llama 2: Open Foundation and Fine-Tuned Chat Models (2023). \texttt{arXiv:2307.09288}
  \item ChatDev: Communicative Agents for Software Development (2023). \texttt{arXiv:2307.07924}
  \item Improving Factuality and Reasoning in Language Models through Multiagent
  Debate (2023). \texttt{arXiv:2305.14325}
  \item Language Models can Solve Computer Tasks (2023). \texttt{arXiv:2303.17491}
  \item AgentBench: Evaluating LLMs as Agents (2023). \texttt{arXiv:2308.03688}
  \item AutoGen: Enabling Next-Gen LLM Applications via Multi-Agent Conversation (2023). \texttt{arXiv:2308.08155}
  \item Describe, Explain, Plan and Select: Interactive Planning with Large
  Language Models Enables Open-World Multi-Task Agents (2023). \texttt{arXiv:2302.01560}
  \item Tool Learning with Foundation Models (2023). \texttt{arXiv:2304.08354}
  \item Tool Learning in the Wild: Empowering Language Models as Automatic Tool
  Agents (2024). \texttt{arXiv:2405.16533}
  \item Exploring Large Language Model based Intelligent Agents: Definitions,
  Methods, and Prospects (2024). \texttt{arXiv:2401.03428}
  \item Large Language Model based Multi-Agents: A Survey of Progress and
  Challenges (2024). \texttt{arXiv:2402.01680}
  \item Exploring Autonomous Agents through the Lens of Large Language Models: A
  Review (2024). \texttt{arXiv:2404.04442}
  \item Large Language Model-Based Agents for Software Engineering: A Survey (2024). \texttt{arXiv:2409.02977}
  \item Creating an LLM-based AI-agent: A high-level methodology towards
  enhancing LLMs with APIs (2024). \texttt{arXiv:2412.13233}
  \item Small LLMs Are Weak Tool Learners: A Multi-LLM Agent (2024). \texttt{arXiv:2401.07324}
  \item LLM With Tools: A Survey (2024). \texttt{arXiv:2409.18807}
  \item Learning to Use Tools via Cooperative and Interactive Agents (2024). \texttt{arXiv:2403.03031}
  \item From Language to Action: A Review of Large Language Models as Autonomous Agents and Tool Users (2025). \texttt{arXiv:2508.17281}
  \item Large Language Model Agent: A Survey on Methodology, Applications and
  Challenges (2025). \texttt{arXiv:2503.21460}
  \item Fundamentals of Building Autonomous LLM Agents (2025). \texttt{arXiv:2510.09244}
  \item From LLM Reasoning to Autonomous AI Agents: A Comprehensive Review (2025). \texttt{arXiv:2504.19678}
  \item AutoTool: Efficient Tool Selection for Large Language Model Agents (2025). \texttt{arXiv:2511.14650}
  \item A Survey of Large Language Model Agents for Question Answering (2025). \texttt{arXiv:2503.19213}
  \item Agentic Reasoning: A Streamlined Framework for Enhancing LLM Reasoning with Agentic Tools (2025). \texttt{arXiv:2502.04644}
  \item Agentic Reasoning and Tool Integration for LLMs via Reinforcement
  Learning (2025). \texttt{arXiv:2505.01441}
  \item Generalizability of Large Language Model-Based Agents: A Comprehensive Survey (2025). \texttt{arXiv:2509.16330}
  \item LLM-based Agentic Reasoning Frameworks: A Survey from Methods to Scenarios (2025). \texttt{arXiv:2508.17692}
  \item Agentic Reasoning for Large Language Models (2026). \texttt{arXiv:2601.12538}
  \item Agentic Tool Use in Large Language Models (2026). \texttt{arXiv:2604.00835}
  \item Agentic Artificial Intelligence (AI): Architectures, Taxonomies, and Evaluation of Large Language Model Agents (2026). \texttt{arXiv:2601.12560}
  \item The Path Ahead for Agentic AI: Challenges and Opportunities (2026). \texttt{arXiv:2601.02749}
\end{enumerate}

\end{document}
